
\exercise
Suppose $T \in \L(V)$. Prove that if both $T$ and $-T$ are positive operators, then $T = 0$.


\exercise
Suppose $T \in \L\paren[\big]{\F4}$ is the operator whose matrix (with respect to the standard basis) is
\[
\mat{cccc}{
2 & -1 & 0 & 0 \\
-1 & 2 & -1 & 0\\
0 & -1 & 2 & - 1 \\
0 & 0 & -1 & 2}.
\]
Show that $T$ is an invertible positive operator.


\exercise
Suppose $n$ is a positive integer and $T \in \L\paren[\big]{\F n}$ is the operator whose matrix (with respect to the standard basis) consists of all $1$'s. Show that $T$ is a positive operator.


\exercise
Suppose $n$ is an integer with $n > 1$. Show that there exists an \by{n}{n} matrix $A$ such that all of the entries of $A$ are positive numbers and $A = A^*\1$, but the operator on $\F n$ whose matrix (with respect to the standard basis) equals $A$ is not a positive operator.


\exercise
Suppose $T \in \L(V)$ is self-adjoint. Prove that $T$ is a positive operator if and only if for every orthonormal basis $e_1, \ldots, e_n$ of $V\1$, all entries on the diagonal of $\M\paren[\big]{ T, (e_1, \ldots, e_n) }$ are nonnegative numbers.


\exercise
Prove that the sum of two positive operators on $V$ is a positive operator.
% \exercisecomment{Surprisingly, the result in this exercise does not seem to follow easily from any of the equivalent characterizations of positive operators in (b) through (f) of \ref{197}.}


\exercise
Suppose $S \in \L(V)$ is an invertible positive operator and $T \in \L(V)$ is a positive operator. Prove that $S+T$ is invertible.


\exercise
Suppose $T \in \L(V)$. Prove that $T$ is a positive operator if and only if the pseudoinverse $T^\dagger$ is a positive operator.\index{pseudoinverse}


\exercise
Suppose $T \in \L(V)$ is a positive operator and $S \in \L(W\3, V)$. Prove that $S^* T S$ is a positive operator on $W\3$.


\exercise
Suppose $T$ is a positive operator on $V\1$. Suppose $v, w \in V$ are such that
\[
Tv = w \andd Tw = v.
\]
Prove that $v = w$.


\exercise
Suppose $T$ is a positive operator on $V$ and $U$ is a subspace of $V$ invariant under $T\3$. Prove that $T|_U \in \L(U)$ is a positive operator on $U$.


\exercise
Suppose $T \in \L(V)$ is a positive operator. Prove that $T^k$ is a positive operator for every positive integer~$k$.


\exercise
Suppose $T \in \L(V)$ is self-adjoint and $\al \in \R{}$.
\begin{subtheorem}
\item
Prove that $T - \al I$ is a positive operator if and only if $\al$ is less than or equal to every eigenvalue of $T\3$.
\item
Prove that $\al I - T$ is a positive operator if and only if $\al$ is greater than or equal to every eigenvalue of $T\3$.
\end{subtheorem}
\exercisesubtheoremend


\exercise
Suppose $T$ is a positive operator on $V$ and $v_1, \ldots, v_m \in V\1$. Prove that
\[
\sum_{\js}^m \sum_{k=1}^m \ip{Tv_k, v_j} \ge 0.
\]
\vspace{-14bp}


\exercise
Suppose $T \in \L(V)$ is self-adjoint. Prove that there exist positive operators $A, B \in \L(V)$ such that
\[
T = A - B \andd \sqrt{T^* T} = A + B \andd AB = BA = 0.
\]
\exercisedisplayend


\exercise
Suppose $T$ is a positive operator on $V\1$. Prove that
\[
\ker \sqrt{T} = \ker T \andd \ran \sqrt{T} = \ran T\3.
\]
\exercisedisplayend


\exercise
Suppose that $T \in \L(V)$ is a positive operator. Prove that there exists a polynomial $p$ with real coefficients such that
$\sqrt{T} = p(T)$. \label{2psqrt}


\begin{comment}
\exercise
Suppose $S \in \L(V)$ and suppose also that $T$ is a positive operator on $V\1$. Prove that $S$ and $T$ commute if and only if $S$ and $\sqrt{T}$ commute.\index{commuting operators|(}

\end{comment}

\exercise
Suppose $S$ and $T$ are positive operators on $V\1$. Prove that $ST$ is a positive operator if and only if $S$ and $T$ commute. \label{2ProdPos} \index{commuting operators}


\exercise
Show that the identity operator on $\F{2}$ has infinitely many self-adjoint square roots.


\begin{comment} Now a result in the Cholesky factorization subsection.
\exercise
Prove that a self-adjoint operator $T \in \L(V)$ is a positive invertible operator if and only if $\ip{Tv, v} > 0$ for every nonzero
$v \in V\1$.\label{7posinvert}


\end{comment}

\exercise
Suppose $T \in \L(V)$ and $e_1, \ldots, e_n$ is an orthonormal basis of $V\1$. Prove that $T$ is a positive operator if and only if there exist
$v_1, \ldots, v_n \in V$ such that
\[
\ip{Te_k, e_j} = \ip{v_k, v_j}
\]
for all $j, k = 1, \ldots, n$.\label{Grammian}
\exercisecomment{The numbers\/ $\br[\big]{\ip{Te_k, e_j}} _{j, k = 1, \ldots, n}$ are the entries in the matrix of\/ $T$ with respect to the orthonormal basis\/ $e_1, \ldots, e_n$.}


\exercise
Suppose $n$ is a positive integer. The \by{n}{n} \emph{Hilbert matrix} is the \by{n}{n} matrix whose entry in row $j$, column $k$ is
$\frac{1}{j+k-1}$. Suppose $T \in \L(V)$ is an operator whose matrix with respect to some orthonormal basis of $V$ is the \by{n}{n} Hilbert matrix. Prove that $T$ is a positive invertible operator.\index{Hilbert matrix}
\exercisecomment{Example: The \by{4}{4} Hilbert matrix is
\[
\mat{cccc}{
1 & \frac{1}{2} & \frac{1}{3} & \frac{1}{4} \\[6bp]
 \frac{1}{2} & \frac{1}{3} & \frac{1}{4} & \frac{1}{5}\\[6bp]
\frac{1}{3} & \frac{1}{4} & \frac{1}{5} & \frac{1}{6}\\[6bp]
\frac{1}{4} & \frac{1}{5} & \frac{1}{6} & \frac{1}{7}
}.
\]}
\vspace{-6bp}


\exercise
Suppose $T \in \L(V)$ is a positive operator and $u \in V$ is such that $\norm{u} = 1$ and $\norm{Tu} \ge \norm{Tv}$ for all $v \in V$ with $\norm{v} = 1$. Show that $u$ is an eigenvector of~$T$ corresponding to the largest eigenvalue of $T\3$.


\exercise
For $T \in \L(V)$ and $u, v \in V\1$, define $\ip{u, v}_T$ by
$
\ip{u, v}_T = \ip{Tu, v}$.
\begin{subtheorem}
\item
Suppose $T \in \L(V)$. Prove that $\ip{\cdot, \cdot}_T$ is an inner product on $V$ if and only if $T$ is an invertible positive operator (with respect to the original inner product $\ip{\cdot, \cdot}$).
\item
Prove that every inner product on $V$ is of the form $\ip{\cdot, \cdot}_T$ for some positive invertible operator $T \in \L(V)$.
\end{subtheorem}
\exercisesubtheoremend


\exercise
Suppose $S$ and $T$ are positive operators on $V\1$. Prove that
\[
\ker(S+T) = \ker S \cap \ker T\3.
\]
\exercisedisplayend


\exercise
Let $T$ be the second derivative operator in Exercise \ref{1100}(b) in Section \ref{7001}. Show that $-T$ is a positive operator.


